\textbf{\emph{Abstract} - 
Urban traffic signal control is often dominated by fixed cycles and local decisions that fail to account for fluctuations in demand and optimize network-wide delay. This study developed and evaluated two novel contributions: Pressure-Based Energy Optimization framework and a QAOA Optimization framework. Energy optimization converts signal control into an Ising-style energy minimization problem, where balancing pressure alignment and neighbor coupling are rewarded and switching phases is penalized. QAOA Optimization encodes each binary traffic-signal phase as a qubit and uses a parameterized QAOA circuit to sample complete network-wide signal configurations from the same Ising Hamiltonian optimized by Classical Global, a network-wide version of the Pressure-Based Energy Optimization framework. The proposed controllers were evaluated against six additional classical and learned traffic-control strategies: Fixed-Time, Max-Pressure, SCOOT, PressLight, MPLight, and CoLight. Testing was performed on both $2\times2$ (four-intersection) and $3\times3$ (nine-intersection) grids using two
complementary experiments: (1) 600-second saturated trials and (2) a 24-hour Seattle case study. Performance was evaluated using throughput, travel time, and waiting time.
}

\textbf{
The energy optimization controllers showed the clearest value when compared to the Max-Pressure controller, since Max-Pressure was similar to these controllers just without many of the energy-specific components. This advantage became most visible as network
size and coordination complexity increased. On the $3\times3$ saturated
network at $\alpha=0.7$, all three energy-based controllers exceeded
Max-Pressure in throughput: 944, 940, and a mean of 980 vehicles completed respectively for the Classical Local, Classical Global, and QAOA controllers compared to 925 vehicles completed for Max-Pressure as well as reduced stopped delay with waiting times of
44.21~s, 43.54~s, and 43.75~s, respectively, compared with 47.16~s for
Max-Pressure. This pattern extended across much of the intermediate
$\alpha$ range with mixed-traffic, where one or more energy-based controllers surpassed
Max-Pressure in throughput as traffic interactions propagated through the
larger network. The 24-hour Seattle case study provided a second test of the first contribution under
time-varying traffic demand. On the smaller $2\times2$ network, Max-Pressure remained particularly strong during sparse periods; however, on the larger $3\times3$ network, the energy-based controllers became substantially more competitive during sustained daytime traffic. During the AM
commute, Classical Global and QAOA achieved average travel times of
76.62~s and 74.93~s, respectively, compared with 79.39~s for Max-Pressure.
During midday, Exact Classical Global and QAOA achieved 81.49~s and
82.03~s, compared with 89.33~s for Max-Pressure. Waiting time showed the
same shift: QAOA achieved 24.95~s during midday compared with 28.37~s for
Max-Pressure, while Exact Classical Global reached 26.06~s. These results indicated that the proposed energy framework was most useful when traffic demand was sufficiently high and the
network sufficiently interconnected for switching stability and neighboring
signal interactions to matter, while its advantage narrowed when the control
problem could be solved primarily from an immediate local pressure imbalance.
}

\textbf{
The value of the second contribution, QAOA Optimization, was evaluated through comparison with Classical Global Optimization, which optimized the same network-wide energy objective classically rather than through a QAOA circuit. On the $2\times2$ saturated
network at $\alpha=0.7$, QAOA reduced travel time by 0.98~s
[95\% CI: $-1.67$ to $-0.21$~s], reduced waiting time by 0.35~s
[95\% CI: $-0.52$ to $-0.14$~s], and completed 14 additional vehicles
[95\% CI: 6-22] relative to Classical Global. On the $3\times3$ saturated network at $\alpha=0.7$, QAOA completed 980 vehicles compared with 940 for Classical Global, an improvement of 40 vehicles
[95\% CI: 27-54]. At $\alpha=0.5$, QAOA completed 1,139 vehicles compared with
1,017 under Classical Global, a substantial increase of 122 vehicles
[95\% CI: 106-138], although this additional discharge came with
13.31~s greater travel time and 1.61~s greater waiting time. Additionally, the 24-hour Seattle case study revealed a complementary QAOA advantage focused on
reducing driver delay. Across the full day on the
$2\times2$ network, QAOA reduced average travel time from 55.67~s under
Classical Global to mean of 53.99~s, a reduction of 1.68~s
[95\% CI: $-1.95$ to $-1.38$~s], and reduced waiting time from 20.19~s
to a mean of 18.81~s, a reduction of 1.38~s
[95\% CI: $-1.55$ to $-1.20$~s]. On the $3\times3$ network, QAOA reduced
full-day travel time from 94.45~s to a mean of 92.96~s
[difference $-1.49$~s; 95\% CI: $-1.78$ to $-1.18$~s] and waiting time
from 30.18~s to a mean of 28.39~s
[difference $-1.78$~s; 95\% CI: $-1.95$ to $-1.61$~s].
Together with the saturated results, these findings showed that QAOA's effect was
regime-dependent: its strongest larger-grid saturated advantage appeared in
vehicle completion, whereas under realistic time-varying Seattle demand its
clearest improvement over Exact Classical Global appeared in waiting
time.
}

\textbf{
The relationship between optimization quality and traffic
quality was tested directly using within-condition energy-traffic
correlations. On the $2\times2$ network, smaller optimality gaps were
significantly associated with lower travel time ($r=0.29$, $p<0.001$), lower
waiting time ($r=0.28$, $p<0.001$), and greater throughput
($r=-0.28$, $p<0.001$); however, on the larger $3\times3$ network, this relationship
became weaker (for example, $r=-0.15$, $p=0.032$ for throughput). In addition, QAOA was subject
to additional validation tests including dedicated
hyperparameter experiments that varied the QAOA mixer angle, cost angle,
circuit depth, and measurement-shot count. Increasing circuit depth from
$p=1$ to $p=3$ raised exact-optimum recovery from 27.23\% to 44.35\%, while
increasing shots from 200 to 1,000 raised recovery from 27.63\% to 48.13\%,
demonstrating that the circuit's approximation quality could be systematically
improved even when corresponding traffic changes remained smaller. Finally,
controlled Gaussian angle perturbation, depolarizing noise, and readout error
were introduced to assess robustness under non-ideal QAOA sampling. Across the
tested low-to-moderate perturbations, travel time and waiting time remained
comparatively stable even when optimality gap and recovery probability changed
more substantially, providing evidence that the closed-loop traffic behavior was
not excessively sensitive to modest changes in QAOA solution quality and
supporting the robustness of the proposed QAOA Optimization framework under
the tested synthetic perturbations.
}

\textbf{
Taken together, the experiments supported two distinct but complementary
contributions to adaptive traffic-signal control. First,
Pressure-Based Energy Optimization extended conventional
pressure-based decision making by combining upstream-downstream traffic
pressure with switching stability and explicit neighbor
coupling in a single interpretable objective. Its strongest advantages emerged when traffic demand was sustained and network interaction
was high, particularly on the larger $3\times3$ grid, indicating that the added
energy terms became increasingly valuable when local queue imbalance alone was
insufficient to evaluate the consequences of a signal decision. Second,
QAOA Global Optimization demonstrated that replacing exact
deterministic minimization of this same network-wide objective with stochastic
sampling of near-optimal global configurations could produce meaningful
closed-loop traffic improvements. On the larger saturated network this benefit
appeared most clearly in vehicle completion, while under Seattle-derived
time-varying demand it appeared most clearly in travel-time and especially waiting-time
reductions relative to Exact Classical Global. This study established where each contribution added value and provided a rigorously tested framework for future traffic control in which the optimization could be matched to network size, traffic demand, and the performance objective (e.g., throughput, travel time, or waiting time) being prioritized.
}

\textbf{\textit{Keywords:}} \textit{Quantum Approximate Optimization Algorithm; traffic signal control; pressure-based control; Ising Hamiltonian; energy minimization; microscopic traffic simulation; SUMO; network-wide optimization; urban traffic congestion; quantum annealer}


\section{Introduction}

Urban traffic congestion is one of the most persistent and costly challenges facing modern cities. In the United States, the 2025 Urban Mobility Report had estimated that the average driver lost 63 hours to traffic congestion in 2024, with nationwide congestion costs reaching approximately \$269 billion annually \cite{schrank2025}. Beyond economic damage, inefficient traffic flow elevates greenhouse gas emissions and lowers air quality for urban residents. Cities and towns together now house 81\% of the world's 8.2 billion people, with two-thirds of future population growth projected to occur in cities through 2050 \cite{un2025}. As global urbanization continues to intensify, the need to develop smarter traffic management has never been greater. At the heart of this problem lies signalized intersections, which are major sources of delay and capacity constraint in urban road networks, where signal timing directly influences vehicle delay \cite{varaiya2013}.

For the purposes of this study, three primary failure modes had been considered:
\textbf{starvation}, where high-demand approaches receive insufficient green
time; \textbf{overflow}, where downstream queues spill back into upstream
approaches and waste available green time; and \textbf{congestion}, where long
queues and repeated acceleration delays reduce vehicle discharge efficiency
through the network.

Traditional fixed-time signal systems operate on predetermined cycle lengths derived from historical patterns and remain widely used for their simplicity and predictability. However, the fundamental limitation of fixed-time systems is that they use predetermined timing plans and therefore cannot react directly to unanticipated real-time demand fluctuations. Although coordinated fixed-time systems can synchronize offsets across intersections, their timing plans remain pre-established rather than continuously optimized from observed traffic conditions. Adaptive traffic control systems improve upon fixed-time signals by adjusting phase timing in response to real-time traffic conditions such as queue lengths. Some systems use local vehicle detectors to modify signal timing at individual intersections, while coordinated systems such as SCOOT coordinate multiple intersections along major roadways to improve traffic flow across a larger network \cite{hunt1981}. Max-pressure control had used queue lengths to select signal phases that minimize the imbalance between upstream and downstream vehicle queues at each intersection \cite{varaiya2013}. Reinforcement learning approaches treat signal control as a sequential decision-making problem, enabling intersections to learn signal timing strategies that reduce congestion through repeated interaction with the traffic environment \cite{agrahari2024}. These adaptive approaches substantially expand beyond static timing, but they
differ in the scope and mechanism of coordination they provide. Max-Pressure,
PressLight, and MPLight primarily generate decisions from local or immediately
downstream traffic conditions, whereas systems such as SCOOT coordinate timing across multiple
intersections and CoLight explicitly exchanges learned information between
neighboring signals \cite{hunt1981,wei2019presslight,chen2020mplight,wei2019colight}. Even in coordinated methods, however,
the complete set of network-wide binary phase assignments is not generally
evaluated as a single shared optimization problem at every decision.

Achieving efficient traffic flow at the network level had remained challenging, however. Neighboring intersections are inherently coupled since vehicles outflow from one intersection to all neighboring ones. This means that it is important to also consider downstream queues as well as upstream queues but also means that signal timing decisions cannot be optimized at the local level without potentially sacrificing the performance of the overall network. For the binary two-phase representation used in this study, simultaneous
optimization of $N$ intersections had required consideration of up to $2^N$
network-wide phase configurations. Exhaustive enumeration is therefore practical
for the $2\times2$ and $3\times3$ networks studied here (16 and 512
configurations, respectively), but its cost had grown exponentially with network
size and becomes increasingly unattractive for larger networks \cite{garey1979}. Pressure-based traffic control had provided a natural basis for such a global
formulation because queue imbalance can be represented as a cost associated with
different signal configurations. In the proposed framework, these traffic-state
costs had been encoded in an Ising-style Hamiltonian containing pressure bias,
downstream-congestion information, switching penalties, and
inter-intersection coupling. Inoue et al. formalize this approach through modeling a $50 \times 50$ intersection lattice with binary signal states and demonstrating that minimizing network-wide flow imbalance is equivalent to minimizing an Ising energy function \cite{inoue2021}. Solving this with a D-Wave 2000Q quantum annealer on a $50 \times 50$ grid, they show that global quantum optimization had outperformed conventional local control, particularly when vehicles travel mostly straight. They further show, analytically, strong spatial correlations between neighboring signals under optimal control, which strongly had supported the idea that network-level coordination becomes increasingly necessary as directional flow starts to dominate. Because Inoue et al. used quantum annealing on a much larger macroscopic
lattice, their observed operating regimes had not been assumed to transfer directly
to the gate-based QAOA controller studied here. The present work instead had tested
QAOA within a microscopic closed-loop SUMO environment and separately had examined
the sensitivity of the QAOA solution to circuit angles, depth, shot count, and
synthetic perturbations. More recent work has continued to apply QAOA to Ising-style traffic congestion
formulations \cite{villanueva2025}, and a 2025 systematic review situates this
line of research within the broader landscape of quantum-assisted traffic
signal control, traffic-flow optimization, and routing \cite{mohammed2025}.
However, an important methodological question had remained unresolved: do
reported traffic benefits arise from network-wide optimization itself, from the
particular quantum search procedure used to approximate that objective, or from
the structure of the objective function being optimized. Addressing this
had required a matched classical global optimizer that had minimized the same
Hamiltonian, together with direct evaluation of how energy quality relates to
downstream traffic performance.

Despite these contributions, several important gaps remain:

\begin{enumerate}
\item Prior quantum traffic-control formulations such as the Inoue et al. model have often emphasized
macroscopic or stylized network models, leaving a need for microscopic
closed-loop evaluation in which individual vehicles, queues, travel times,
waiting times, and vehicle completion had been measured directly.

\item Quantum and energy-based traffic controllers should be evaluated against
strong established adaptive alternatives rather than only Fixed-Time or simple
queue-based baselines. Comparisons with Max-Pressure, SCOOT, and modern
reinforcement-learning methods are needed to determine whether any observed
benefit persists against competitive classical control strategies.

\item A matched exact global optimizer is needed to separate the benefit of
network-wide optimization from the benefit of the QAOA search procedure itself.
Without such a benchmark, improved traffic performance cannot be attributed
specifically to QAOA rather than to the transition from local to global
optimization.

\item Evaluation on a single small network cannot establish whether controller
behavior persists as coordination complexity had increased. Comparing $2\times2$
and $3\times3$ networks had provided an initial controlled scaling test while
retaining a tractable exact-global benchmark.

\item Minimizing an Ising objective is only useful if lower energy had been associated
with improved downstream traffic performance. The relationship between
optimized energy, travel time, waiting time, and throughput therefore had required
direct statistical validation rather than being assumed.

\item QAOA traffic-control results can depend on circuit angles, circuit depth,
measurement shots, stochastic sampling, and quantum noise. Hyperparameter
sensitivity, repeated-run uncertainty, and controlled noise perturbations are
therefore necessary to determine whether observed traffic effects are robust.

\item Finally, synthetic saturated traffic alone had not shown whether the same
controller behavior persists under realistic time-varying urban demand.
A Seattle-derived 24-hour case study had provided a complementary test using
empirical hourly demand and directional-flow information.
\end{enumerate}

Throughout this paper, these gaps had been addressed through two linked scientific
contributions. The first is a \textbf{Pressure-Based Energy Optimization}
framework that had represented traffic-signal selection through a shared Ising-style
objective incorporating traffic pressure, downstream congestion, phase-switching
cost, and neighboring-intersection coupling. This framework had been evaluated using
both Classical Local and Exact Classical Global optimization in order to isolate
the effect of optimization scope. The second contribution is a
\textbf{QAOA-Based Global Optimization} controller that had searched the same
network-wide energy landscape through a parameterized quantum circuit and
stochastic measurement sampling.

Rather than treating QAOA as a standalone replacement for classical traffic
control, the study had evaluated three separate questions: whether the proposed
energy formulation had improved traffic control relative to established adaptive
baselines, whether global energy optimization had improved on local optimization,
and whether QAOA had produced useful downstream traffic behavior relative to an
exact matched-objective global optimizer. The study further had tested whether
optimization quality had predicted traffic quality and how these relationships change
with network size, traffic regime, QAOA configuration, and synthetic noise.

Classical Local Optimization had applied the same energy formulation independently
at each signal, whereas Exact Classical Global Optimization exhaustively
had evaluated the complete network configuration. Their comparison had isolated the
effect of network-wide coordination. Exact Classical Global also serves as a
matched-objective benchmark for QAOA because it had minimized the same Ising
Hamiltonian exactly; differences between the two global controllers therefore
reflect QAOA's approximate stochastic search and the resulting closed-loop
traffic trajectory rather than merely the use of a global objective.

Because the Ising Hamiltonian is a surrogate control objective rather than a
direct measure of future vehicle delay, the study had not assume that lower
energy necessarily had produced better traffic performance. Instead, this
relationship had been tested explicitly by comparing QAOA optimality gap with
downstream travel time, waiting time, and throughput within matched traffic
conditions.

The study also had examined whether the observed QAOA behavior is robust to its
own stochasticity and implementation choices by evaluating repeated-run
variability, circuit-angle sensitivity, circuit depth, measurement-shot count,
and three controlled synthetic perturbation mechanisms.

\begin{itemize}
\item \textbf{RQ1:} Under which traffic regimes and network sizes does
Pressure-Based Energy Optimization provide advantages relative to established
adaptive traffic-signal controllers?

\item \textbf{RQ2:} When the same network-wide Ising objective had been optimized,
how does QAOA compare with Exact Classical Global Optimization in optimization
quality, downstream traffic performance, statistical variability, and
sensitivity to circuit configuration and synthetic noise?

\item \textbf{Hypothesis:} Pressure-Based Energy Optimization had been expected to
be most competitive under moderate-to-high demand and larger-network conditions
where downstream congestion, switching stability, and neighboring-intersection
coordination become important. Exact Classical Global Optimization had been expected
to provide greater benefit over Classical Local Optimization when traffic
movements create stronger interactions across intersections. QAOA Global
Optimization had been expected to produce traffic outcomes that differ from Exact
Classical Global Optimization because its stochastic search can select
near-optimal global configurations rather than repeatedly selecting the exact
instantaneous energy minimum. These differences had been expected to be
regime-dependent rather than universally favorable: under sparse demand,
highly predictable directional flow, or traffic dominated by a clear local
pressure imbalance, Max-Pressure, cycle-based control, or learned adaptive
controllers may remain competitive or outperform the energy-based approaches.
As network size had increased, the QAOA search had been expected to face a more difficult
optimization landscape, making both exact-optimum recovery and the relationship
between energy quality and downstream traffic performance important empirical
questions.

\item \textbf{Aim:} This study had aimed to comparatively evaluate nine
traffic-signal controllers---Fixed-Time, Max-Pressure, SCOOT, PressLight,
MPLight, CoLight, Classical Local Optimization, Exact Classical Global
Optimization, and QAOA Global Optimization---in SUMO microsimulations of both
$2\times2$ and $3\times3$ urban grids. Two complementary traffic experiments
are used: (1) 600-second saturated trials in which $\alpha$ had been varied from
0.0 to 1.0 to characterize controller behavior across turning regimes and
network sizes; and (2) a 24-hour Seattle-derived case study using empirical
hourly demand and directional-flow data.

Traffic performance had been evaluated using completed-vehicle throughput, average
travel time, and average waiting time. QAOA had additionally been assessed through
repeated-run variability, bootstrap confidence intervals, and effect-size
comparisons with Exact Classical Global Optimization. Exact enumeration is used
to quantify optimality gap and exact-optimum recovery, correlation analysis is
used to test whether improved Ising optimization had predicted better downstream
traffic, and dedicated sensitivity experiments examine circuit parameters,
depth, shot count, and synthetic quantum perturbations.
\end{itemize}
